\chapter{Wykład VII'}

\section{Grupy rozwiązalne. Komutator i Komutant} 

\begin{context}
	Wkraczamy w teorię, która historycznie dała nazwę całej dziedzinie "rozwiązywania równań". Zbadamy, jak bardzo grupa różni się od grupy przemiennej. Narzędziem do tego pomiaru jest komutator.
\end{context}

\subsection{Komutatory i ich własności}

Niech $G$ będzie dowolną grupą.

\begin{definition}[Komutator]
	Dla elementów $x,y \in G$ definiujemy ich komutator jako element:
	\[
		\left[ x,y \right] := xyx^{-1}y^{-1}.
	\] 
\end{definition} 

\begin{remark}[Własności komutatora]
	Komutator "mierzy" nieprzemienność elementów:
	\begin{itemize}
		\item $xy=yx \iff xyx^{-1}y^{-1} = e \iff \left[ x,y \right] = e$.
		\item Grupa $G$ jest abelowa $\iff \forall\left( x,y \in G \right) \left( \left[ x,y \right] = e \right)$.
		\item Zachowanie przy homomorfizmach: Niech $\varphi : G \to H$ będzie homomorfizmem. Wtedy:
		\[ \varphi \left( \left[ x,y \right] \right) = \varphi(xyx^{-1}y^{-1}) = \varphi(x)\varphi(y)\varphi(x)^{-1}\varphi(y)^{-1} = \left[ \varphi (x), \varphi \left( y \right) \right]. \]
		\item Odwracanie: $[x,y]^{-1} = (xyx^{-1}y^{-1})^{-1} = yxy^{-1}x^{-1} = [y,x]$.
	\end{itemize}
\end{remark}

\begin{definition}[Komutant podgrup]
	Niech $H, T \le G$. Wzajemnym komutantem tych podgrup nazywamy podgrupę generowaną przez wszystkie możliwe komutatory ich elementów:
	\[ \left[ H,T \right] := \langle \{\left[ x,y \right]: x\in H, y\in T \}  \rangle. \]
	W szczególności, jeśli $H=T=G$, to podgrupę $[G,G]$ nazywamy \textbf{komutantem grupy} $G$ i oznaczamy $G'$ lub $G^{(1)}$.
\end{definition} 

\begin{fact}
	Jeżeli $H\unlhd G$ oraz $T\unlhd G$, to ich komutant też jest podgrupą normalną:
	\[ \left[ H,T \right] \unlhd G. \]
	W szczególności komutant grupy $G' = [G,G]$ jest zawsze podgrupą normalną w $G$.
\end{fact}
\begin{proof}
	Pokażemy, że sprzężenie elementu z $[H,T]$ nadal należy do $[H,T]$. Wystarczy sprawdzić to na generatorach (czystych komutatorach).
	Niech $h \in H, t \in T, g \in G$. Rozważmy sprzężenie komutatora $[h,t]$ przez $g$:
	\begin{align*}
		g [h,t] g^{-1} &= g (hth^{-1}t^{-1}) g^{-1} \\
		&= (ghg^{-1}) (gtg^{-1}) (gh^{-1}g^{-1}) (gt^{-1}g^{-1}) \\
		&= (ghg^{-1}) (gtg^{-1}) (ghg^{-1})^{-1} (gtg^{-1})^{-1} \\
		&= [ghg^{-1}, gtg^{-1}].
	\end{align*}
	Ponieważ $H$ i $T$ są normalne, to $h' = ghg^{-1} \in H$ oraz $t' = gtg^{-1} \in T$.
	Zatem wynik jest komutatorem $[h', t']$, który z definicji należy do $[H,T]$.
\end{proof} 

\subsection{Szereg pochodny i Abelianizacja}

Tworzymy ciąg podgrup, biorąc kolejne komutanty komutantów.

\begin{definition}[Szereg pochodny]
	Definiujemy indukcyjnie ciąg podgrup (zwany szeregiem pochodnym):
	\begin{itemize}
		\item $G^{\left( 0 \right) }:= G$,
		\item $G^{\left( 1 \right) }= G' := \left[ G,G \right] $ (komutant),
		\item $G^{(n+1)} := \left( G^{(n)} \right)' = \left[ G^{\left( n \right) }, G^{\left( n \right) } \right] $.
	\end{itemize}
	Zauważmy, że tworzy to ciąg podgrup normalnych:
	\[ G = G^{(0)} \unrhd G^{(1)} \unrhd G^{(2)} \unrhd \ldots \]
\end{definition} 

Komutant $G'$ jest najmniejszą podgrupą normalną, po której "zababicie" (iloraz) daje grupę przemienną.

\begin{definition}[Abelianizacja]
	Grupę ilorazową:
	\[ G^{ab} := \mfaktor{G}{G'} \]
	nazywamy \textbf{abelianizacją} grupy $G$. Jest to największa abelowa grupa ilorazowa $G$.
\end{definition} 

\begin{corollary}
	$G^{ab}$ jest zawsze grupą abelową.
\end{corollary}
\begin{proof}
	Niech $\bar{x}, \bar{y} \in G/G'$. Wtedy $[\bar{x}, \bar{y}] = \overline{[x,y]}$. Ale $[x,y] \in G'$, więc w ilorazie jest elementem neutralnym. Zatem $\bar{x}\bar{y}\bar{x}^{-1}\bar{y}^{-1} = e$, czyli $\bar{x}\bar{y} = \bar{y}\bar{x}$.
\end{proof}

\begin{theorem}[O uniwersalności komutanta]
	Niech $N\unlhd G $. Wtedy:
	\[ \mfaktor{G}{N} \text{ jest abelowa} \iff G' \subseteq N. \]
\end{theorem} 
\begin{proof}
	$(\Rightarrow)$ Załóżmy, że $G/N$ jest abelowa.
	Weźmy dowolne $x,y \in G$. W grupie ilorazowej mamy $(xN)(yN) = (yN)(xN)$.
	To oznacza $xyN = yxN$, czyli mnożąc stronami: $x^{-1}y^{-1}xyN = N$.
	Ale $x^{-1}y^{-1}xy = [x^{-1}, y^{-1}]$, a zbiór takich elementów generuje $G'$. 
	Zatem każdy komutant należy do $N$, więc $G' \subseteq N$.
	
	$(\Leftarrow)$ Załóżmy, że $G' \subseteq N$.
	Wtedy dla dowolnych $x,y$ mamy $[x,y] \in N$.
	To oznacza, że $xyx^{-1}y^{-1} \in N$, czyli w ilorazie $xyN = yxN$. Grupa jest abelowa.
\end{proof} 

\subsection{Przykłady komutantów}

Obliczmy komutanty $G'$ dla znanych grup.

\begin{enumerate}[label=\arabic*)]
	\item $S'_{3} = A_3$.
	\begin{proposition}
		$S_3$ nie jest abelowa, więc $S_3' \neq \{e\}$. $S_3/A_3 \cong \mathbb{Z}_2$ (abelowa), więc $S_3' \subseteq A_3$. Jedyną nietrywialną podgrupą $A_3$ jest samo $A_3$.
		Konkretnie: $[(1,2), (1,3)] = (1,2)(1,3)(1,2)(1,3) = (1,3,2)(1,3) = (1,2,3) \in A_3$.
	\end{proposition}

	\item $A'_{4} = V_4 = \{\id, (1,2)(3,4), (1,3)(2,4), (1,4)(2,3)\} $.
	\begin{proposition}
		$A_4$ ma podgrupę normalną $V_4$. Iloraz $A_4/V_4 \cong \mathbb{Z}_3$ (abelowa), zatem $A_4' \subseteq V_4$.
		Z drugiej strony obliczamy komutator cykli: $[(1,2,3), (1,2,4)] = (1,2)(3,4) \in V_4$.
	\end{proposition}

	\item $\mathcal{Q}'_{8} = \{\mathcal{I}, -\mathcal{I}\} = Z(\mathcal{Q}_8)$.
	\begin{proposition}
		Mamy $ij = k$, ale $ji = -k$. Zatem $[i,j] = iji^{-1}j^{-1} = k(-i)(-j) = k \cdot k = -1$.
		Iloraz $Q_8 / \{1, -1\}$ to grupa Kleina $V_4$ (abelowa), więc komutant to dokładnie $\{-1, 1\}$.
	\end{proposition}

	\item $\forall\left( n \right)\left( S'_{n} = A_{n} \right)$.
	\begin{proposition}
		Dla $n \ge 2$ znak komutatora to $\sgn(xyx^{-1}y^{-1}) = \sgn(x)\sgn(y)\sgn(x)\sgn(y) = 1$, więc $S_n' \subseteq A_n$.
	\end{proposition}

	\item $\forall\left( n \ge 5 \right) \left( A'_{n} = A_{n} \right)$.
	\begin{proposition}
		To kluczowy fakt! Grupy $A_n$ dla $n \ge 5$ są \textbf{grupami prostymi} (nie mają właściwych podgrup normalnych) i są nieprzemienne.
		Skoro $A_n' \unlhd A_n$, to $A_n'$ musi być równe $A_n$ (bo nie może być trywialne).
		Taka grupa nazywa się \textbf{doskonała} (perfect).
	\end{proposition}
\end{enumerate}

\section{Rozwiązalność}

Idea rozwiązalności polega na sprawdzaniu, czy biorąc kolejne komutanty, dojdziemy w końcu do grupy trywialnej $\{e\}$. Jeśli tak - grupa jest "niezbyt skomplikowana" (złożona z warstw abelowych).

\begin{definition}[Grupa rozwiązalna]
	Grupę $G$ nazywamy \textbf{rozwiązalną}, gdy jej szereg pochodny kończy się na $\{e\}$.
	Tzn. istnieje $n\in \N$, takie że:
	\[ G^{(n)} = \{e\}. \]
	Najmniejsze takie $n$ nazywamy stopniem rozwiązalności (lub długością pochodną).
\end{definition} 

\begin{remark}
	$G$ jest abelowa $\iff G' = \{e\} \iff G$ jest rozwiązalna stopnia 1.
\end{remark} 

\textbf{Analiza rozwiązalności dla przykładów (zastosowanie definicji)}

Sprawdźmy ciągi $G \unrhd G^{(1)} \unrhd G^{(2)} \dots$ dla powyższych grup.

\begin{enumerate}[label=\arabic*)]
	\item $G = S_3$.
	\[ S_3 \triangleright S_3^{(1)} = A_3 \triangleright S_3^{(2)} = A_3' = \{e\}. \]
	(Bo $A_3 \cong \mathbb{Z}_3$ jest abelowa, więc jej komutant to $\{e\}$).
	\textbf{Wniosek:} $S_3$ jest rozwiązalna (stopnia 2).

	\item $G = S_4$.
	\[ S_4 \triangleright S_4^{(1)} = A_4 \triangleright S_4^{(2)} = V_4 \triangleright S_4^{(3)} = \{e\}. \]
	(Bo $V_4 \cong \mathbb{Z}_2 \oplus \mathbb{Z}_2$ jest abelowa).
	\textbf{Wniosek:} $S_4$ jest rozwiązalna (stopnia 3).

	\item $G = \mathcal{Q}_8$.
	\[ \mathcal{Q}_8 \triangleright \mathcal{Q}_8^{(1)} = \{1, -1\} \triangleright \mathcal{Q}_8^{(2)} = \{1\}. \]
	(Bo $\{1, -1\}$ jest grupą cykliczną rzędu 2, więc abelową).
	\textbf{Wniosek:} $\mathcal{Q}_8$ jest rozwiązalna.

	\item $G = A_5$ (i ogólnie $A_n$ dla $n \ge 5$).
	\[ A_5 \triangleright A_5^{(1)} = A_5 \triangleright A_5^{(2)} = A_5 \dots \]
	Ciąg się ustabilizował i nigdy nie osiągnie $\{e\}$.
	\textbf{Wniosek:} $A_5$ (oraz $S_n$ dla $n \ge 5$) \textbf{NIE} są rozwiązalne.
	
	\textit{Komentarz historyczny:} To właśnie ten fakt w teorii Galois odpowiada za brak ogólnych wzorów na pierwiastki wielomianów stopnia $\ge 5$.
\end{enumerate}
